\section{What the positive cyclic-trace decomposition does not imply}
\label{sec:trace-index-obstruction}

The identity $k_n=b_n+3c_n$, with $c_n\ge0$, is stronger than a
coefficient congruence. The following countermodel shows that this
strengthening, even together with several numerical properties of the
two actual quotients, does not transfer large strict indices. The series
constructed here are \emph{formal rational characters}. They are not
asserted to be invariant-ring Hilbert series, and the series denoted $G$
is not the actual auxiliary series $B_n$. Accordingly the result does
not settle either strict index of the tensor family.

\begin{theorem}[PNEW-014: a positive trace can mask large strict indices]
For every even $n\ge10$, put
\[
 A=n^2,\qquad D=n^3-3n^2+3,\qquad E=2\binom n3+1,
 \qquad S=\lceil\sqrt n\rceil,\qquad M=2^{n-1}.
 \tag{PNEW-Q1}
\]
There are proper rational functions $F,G\in\mathbb Q(t)$ with globally
quasipolynomial, nonnegative integral coefficients $f(s),g(s)$ such that:
\begin{enumerate}
\item $f(0)=g(0)=1$, and $f(s)=g(s)=0$ for $1\le s<S$.
\item $f=g+3c$ for a nonnegative integral sequence $c$. This is the
dimension--trace identity of an actual rational $C_3$ representation
in each degree.
\item Their pole orders at one are $D$ and $E$, respectively, and
\[
 F(t^{-1})=-t^A F(t),\qquad G(t^{-1})=-t^A G(t).
 \tag{PNEW-Q2}
\]
\item For every $s\ge0$ both coefficient sequences satisfy the tensor
ambient bound
\[
        0\le g(s)\le f(s)\le\binom{ns+n^3-1}{n^3-1}.
 \tag{PNEW-Q3}
\]
\item $G$ has a simple primitive-$M$ pole with residue a unit at
every prime above three, and no primitive-$M3^j$ pole for $j\ge1$.
The function $F$ has no primitive-$M$ pole and has a simple
primitive-$3M$ pole.
\item In the strict constituent conventions,
\[
 s(g)\ge 2^{n-5},\qquad p(g)\ge s(g),\qquad
 s(f)=S-1,\qquad p(f)\le n^3-4n^2+4.
 \tag{PNEW-Q4}
\]
\end{enumerate}
\end{theorem}
\begin{proof}
Write
\[
 \begin{aligned}
 L&=M/2,\qquad r=(L-A-2)/2,\qquad a=\lfloor r/2\rfloor,\\
 d_E&=E+1-A,\qquad d_D=D+1-A,\qquad K=D-A-2S.
 \end{aligned}
 \tag{PNEW-Q5}
\]
These are positive integers. The elementary inequalities needed below are
\[
 E\ge3,\quad E+1\le D,\quad S<d_E,\quad
 0<K<D,\quad A+2(S-1)<L,\quad r\ge2S,
 \quad a\ge L/8=2^{n-5}.
 \tag{PNEW-Q6}
\]
For example $L\ge2A+4$ holds at $n=10$ and is preserved on replacing
$n$ by $n+2$, since $L$ is multiplied by four. This proves the last
inequality and also $r\ge2S$. The other inequalities follow directly
from $E=\frac13n(n-1)(n-2)+1$, $D=n^3-3n^2+3$, and $S\le n$;
their least endpoint is $n=10$.

Define the polynomial
\[
 q(X)=\binom{X+E}{E}-\binom{X+A-1}{E}
 \quad\hbox{and}\quad
 Q(t)=\frac{1-t^{d_E}}{(1-t)^{E+1}}.
 \tag{PNEW-Q7}
\]
Binomial expressions here denote their polynomial extensions.
For every nonnegative integer $s$, the coefficient of $t^s$ in $Q$
is $q(s)$. Before $s=d_E$ the second binomial vanishes, since its
upper argument is a nonnegative integer less than $E$.
Equivalently
\[
 q(s)=\sum_{i=0}^{d_E-1}\binom{s-i+E-1}{E-1}_{+},
 \tag{PNEW-Q8}
\]
where the summand is zero for $s<i$. This proves that $q(s)$ is
positive and nondecreasing on nonnegative integers. Its polynomial
degree is $E-1$ with positive leading coefficient.

Let $C$ be the least integer at least $q(a)$ congruent to one modulo
three, and set $h=(C-1)/3$. Thus
\[
                  q(a)\le C\le q(a)+2,\qquad h\ge0.
 \tag{PNEW-Q9}
\]
The common correction, positive polynomial character, and two waves are
\[
\begin{split}
 J(t)&=-\sum_{j=1}^{S-1}\frac{q(j)t^j}{1-t^{A+2j}},\\
 Q_0(t)&=\frac{t^S(1+t)^K}{(1-t)^D},\\
 T_G(t)&=\frac{t^r(1-t^2)}{1+t^L},\qquad
 T_F(t)=\frac{t^r(1-t^2)(1+t^L)}{1-t^L+t^{2L}},\\
 G(t)&=Q(t)+J(t)+C T_G(t),\\
 F(t)&=Q(t)+J(t)+3Q_0(t)+T_F(t).
\end{split}
 \tag{PNEW-Q10}
\]
All are proper rational functions with root-of-unity poles. The coefficients
of $J,T_G,T_F$ are globally periodic: all their displayed numerators
have degree less than the relevant period denominator after rationalizing
to $1-t^P$. The coefficients of $Q_0$ are the global polynomial
\[
 q_0(X)=\sum_{j=0}^K\binom Kj
                \binom{X-S-j+D-1}{D-1}.
 \tag{PNEW-Q11}
\]
For $0\le s<S+j$, the upper argument in each term is nonnegative
and less than $D-1$, because $S+K<D$. Thus this polynomial agrees
with the power-series coefficients even before each summand starts.
Consequently both $F$ and $G$ are globally quasipolynomial.

The coefficients of $T_G$ have absolute value at most one; those of
$T_F$ have absolute value at most two. Indeed the coefficient sequences
in $z=t^L$ of $1/(1+z)$ and $(1+z)/(1-z+z^2)$ are respectively
$(1,-1)$ and $(1,2,1,-1,-2,-1)$ periodically. The two shifts separated
by two do not overlap because $L>2$. Both waves start at degree $r$.

Put $B=\sum_{j=1}^{S-1}q(j)$. Since $S<d_E$, the hockey-stick identity
gives the exact estimate
\[
 B=\binom{E+S}{E+1}-1
   =\frac{S}{E+1}q(S)-1,
 \qquad q(S)-B-2>0.
 \tag{PNEW-Q12}
\]
The second inequality follows from $S<E$, $q(S)\ge E+1$.
In degrees $1\le s<S$, precisely the $j=s$ summand of $J$
cancels $q(s)$; all waves and $Q_0$ vanish. The constant terms are one.
For $S\le s<r$, $q(s)-B>0$ proves $g(s)>0$.
For $s\ge r$, use $r\ge2a$ to obtain
\[
                         q(r)\ge2q(a).
 \tag{PNEW-Q13}
\]
To verify this term by term in (PNEW-Q8), put $u=a-i\ge1$.
The first two factors in
$\binom{2u+E-1}{E-1}/\binom{u+E-1}{E-1}$ already have product
$2(2u+1)/(u+2)\ge2$. If $u=0$, the corresponding numerator at
$r-i$ is at least $E\ge3$; terms absent at $a$ cause no problem.
Equations (PNEW-Q9), (PNEW-Q12), and monotonicity now give
\[
               g(s)\ge q(r)-B-C\ge q(a)-B-2>0.
 \tag{PNEW-Q14}
\]
Thus $g$ is nonnegative and integral.

The elementary rational identity
\[
 T_F-T_G=\frac{3t^{r+L}(1-t^2)}{1+t^{3L}}=:3T_0
 \tag{PNEW-Q15}
\]
gives
\[
                         F-G=3(Q_0+T_0-hT_G).
 \tag{PNEW-Q16}
\]
The coefficients of $T_0$ have absolute value at most one and start
after degree $r$. Before degree $r$ the right-hand coefficient is
therefore nonnegative. After that degree, the positive coefficients
of $Q_0$ are nondecreasing and satisfy
\[
 q_0(r)\ge\binom{r-S+D-1}{D-1}
          \ge\binom{a+E}{E}\ge q(a)\ge h+1.
 \tag{PNEW-Q17}
\]
Here $r-S\ge a$, $D-1\ge E$, and $q(a)\ge2$ justify all steps.
This proves $c(s)\ge0$ and hence $f(s)\ge g(s)$. The representation
\[
       \mathbb Q^{\,g(s)+c(s)}\ \oplus\
                 \mathbb Q(\omega)^{\,c(s)}
 \tag{PNEW-Q18}
\]
has dimension $f(s)$ and cyclic trace $g(s)$; here a generator of
$C_3$ acts by multiplication by $\omega$ on each quadratic field
summand. No multiplication between degrees is asserted.

The pole at one of $Q$ has order $E$. The pole of $Q_0$ has order
$D>E$, while $J$ has at most simple poles and both waves are regular
at one. This proves the two dimensions. Direct substitution $t\mapsto
t^{-1}$ in each term of (PNEW-Q10) gives the common multiplier
$-t^A$: use $E+1-d_E=A$, $D-2S-K=A$, and $L-2-2r=A$.
Thus (PNEW-Q2) is an exact rational functional equation.

We give an injection proving the ambient bound, to avoid any asymptotic
estimate. For $s\ge S$,
\[
              f(s)\le\binom{s+E}{E}+3q_0(s)+2.
 \tag{PNEW-Q19}
\]
Let $N=n^3$ and use variables $y_1,\ldots,y_N$, with
$D+1<N$. A monomial counted by $q_0(s)$ is a tuple of free exponents
$u_1,\ldots,u_D\ge0$ and bits $e_1,\ldots,e_K\in\{0,1\}$ with
$\sum u_i+\sum e_i=s-S$. In copy $b\in\{1,2,3\}$ encode it by
$y_i^{2u_i+e_i}$ for $i\le K$, $y_i^{2u_i}$ otherwise,
and $y_{D+1}^b$. Pad with $y_N$ to total degree $ns$.
The unpadded degree is at most $2(s-S)+3<ns$.
Parity recovers every bit and exponent, so this is injective.
Encode a monomial of degree $s$ in $E+1$ variables by doubled
exponents in the first $E+1$ variables and marker $y_{D+1}^4$;
its unpadded degree $2s+4$ is at most $ns$.
The two final objects use markers $y_{D+1}^5,y_{D+1}^6$ and
zero free exponents. Again pad with $y_N$.
These five marker types give pairwise disjoint images in degree $ns$.
This proves (PNEW-Q3) for $s\ge S$; the initial degrees are immediate.

All poles of $J$ have orders at most $A+2(S-1)<L$.
At a primitive $M=2L$ root $\zeta$, the unique pole in $G$ comes from
$C T_G$, and
\[
       \lim_{t\to\zeta}(1-t^M)G(t)
                   =2C\zeta^r(1-\zeta^2).
 \tag{PNEW-Q20}
\]
This is a unit at every prime above three: $C\equiv1\pmod3$,
and reduction preserves the exact dyadic order of $\zeta$.
There are no primitive-$M3^j$ poles in $G$ for $j\ge1$.
The only denominator of $F$ with two-part at least $M$ is
$1-t^L+t^{2L}=\Phi_{3M}(t)$. Its numerator in $T_F$ is nonzero at
every such root. Thus $F$ has a simple primitive-$3M$ pole and
no primitive-$M$ pole.

It remains to verify the strict indices, including the off-residue tests.
Use any common period of the periodic terms; changing to the least
period merely removes duplicate constituents. The constituent of $G$
whose $T_G$ coefficient is $-1$ has value at $X=a$ at most
$q(a)-C\le0$, since the $J$ contribution is nonpositive.
It is nonzero because its leading degree is $E-1$.
This gives $s(g)\ge a\ge2^{n-5}$ and $p(g)\ge s(g)$.

Every constituent of $F$, evaluated at any integer $X\ge S$, is
at least $q(X)-B-2+3q_0(X)>0$ by (PNEW-Q12).
Consequently $s(f)\le S-1$. Its actual zero at $s=S-1$, in a
nonzero constituent of degree $D-1$, gives the reverse inequality.
Finally let $u=d_D$. Since $u\ge d_E$, the polynomial $q(X+u)$ is
the difference of two products of $E$ linear factors with positive
constant terms; each first constant exceeds its corresponding second
constant by $d_E$. Every coefficient below the cancelling top degree
is strictly positive. Its constant term exceeds $B+2$ by monotonicity
and (PNEW-Q12). Each term of $q_0(X+u)$ is a product of $D-1$
linear factors with positive constants, because $u>S+K$.
The periodic terms affect only the constant coefficient of each
constituent. Therefore every constituent of $F(X+u)$ has nonnegative
coefficients, proving $p(f)\le u=D+1-A$.
\end{proof}

The two waves in (PNEW-Q10) begin with the exact-order warning example
of PNEW-013; the dimension and duality construction was proposed by
the Kronecker branch. The positive mask, auxiliary-dimension match,
exponentially large auxiliary indices, initial vanishing range, and
ambient injection above make the scope of the obstruction precise.
They show that those structural data do not imply a strict-index
transfer. An argument using further coefficients or algebraic structure
of the \emph{actual} $B_n$ and tensor ring remains possible and is needed.
